\documentclass[11pt,a4paper]{article}
\usepackage[T1]{fontenc}
\usepackage{geometry}
\usepackage{amsmath}
\usepackage{amssymb}
\usepackage{booktabs}
\usepackage{graphicx}
\usepackage{xcolor}
\usepackage{array}
\usepackage{mathptmx}
\usepackage{listings}
\usepackage{hyperref}

\geometry{margin=1in}

\definecolor{codebg}{gray}{0.95}
\lstset{
  basicstyle=\ttfamily\small,
  backgroundcolor=\color{codebg},
  breaklines=true,
  frame=single,
  columns=fullflexible,
  showstringspaces=false
}

\title{Running an LU-Factorization GPU Solver under Nixnan\\
\large Exp-Binade Histogramming, an Exponential \texttt{SAMPLING} Sweep, and Two Tool-Level Findings}
\author{SC26 Tutorial / Nixnan Experiments}
\date{}

\begin{document}
\maketitle

\section{Executive Summary}

This report documents \texttt{lu\_solve\_gpu}, a self-contained CUDA port of
the LU-factorization solver from the FPChecker tutorial
(\texttt{tutorial/example\_1/lu\_solve.cpp} in the sibling
\texttt{SC26-Tutorial-NixNan} repository), run under nixnan
(\url{https://github.com/parfloat/nixnan}) on an NVIDIA RTX 3090
(compute capability 8.6). The experiment has two parts:

\begin{enumerate}
  \item \textbf{Exp-binade histogram gathering}: an exponent-range
    (\emph{binade}) monitoring configuration covering the one floating-point
    format the solver actually uses (\texttt{f64}), bucketed at a resolution
    of 4 exponent values per bin across the full valid double exponent
    range, reproduces the tutorial's central story on the GPU: the
    well-conditioned matrix solves cleanly (residual $\approx 10^{-16}$,
    zero exceptions) while the ill-conditioned matrix collapses to NaN
    (316 NaN / 4 Infinity / 4 Division-by-zero events in FP64 operations).
  \item \textbf{An exponential \texttt{SAMPLING} sweep} (every~1, then
    every~2, every~4, every~8 kernel invocation, 16 runs per level, then
    reset back to every~1) was implemented and run as requested. In the
    process we found that nixnan's \texttt{SAMPLING} feature has
    \textbf{no observable effect in this build} due to a key-mismatch bug
    between the launch counter and the sampling lookup (Section
    \ref{sec:sampling-bug}), and that \texttt{BIN\_SPEC\_FILE}'s documented
    \texttt{doublings} field is not read anywhere in the current
    \texttt{src/fp-histogram.cu} (Section \ref{sec:doublings}). Both are
    tool-level observations, not properties of the solver being tested.
\end{enumerate}

All configuration files, the CUDA source, and the run logs referenced here
live alongside this report in \texttt{SC26-expts/lu\_solve\_gpu/}.

\section{Provenance and Algorithm}

\texttt{lu\_solve\_gpu.cu} solves $Ax = \mathbf{1}$ for a $20\times20$ matrix
$A$ loaded from CSV, using LU factorization with partial pivoting
($PA = LU$) followed by forward/backward substitution -- the same algorithm,
including the same ``snap near-machine-epsilon values to exactly zero''
elimination rule, as the CPU version it was ported from. Matrices are
flattened into row-major device buffers; the pivot search, row swap, and
elimination steps run as CUDA kernels (Table~\ref{tab:kernels}), while the
two substitution passes run as single-thread kernels because of their
sequential row-to-row data dependency. Two input matrices are used
throughout:

\begin{itemize}
  \item \texttt{matrix.csv} -- well-conditioned; solves to residual
    $\approx 9.68\times10^{-16}$.
  \item \texttt{bad\_matrix.csv} -- deliberately ill-conditioned (several
    entries perturbed by $\sim10^{-6}$ relative to an exact tridiagonal
    structure). A pivot in the elimination step is rounded to exactly
    \texttt{0.0} by the epsilon-clamp rule, the subsequent
    backward-substitution division by that zero pivot produces
    $\pm\infty$, and the resulting NaNs propagate through the rest of the
    solution and the residual.
\end{itemize}

\begin{table}[h]
\centering
\begin{tabular}{ll}
\toprule
Kernel & Role \\
\midrule
\texttt{init\_identity\_kernel} & Initialize $L$, $P$ as identity matrices \\
\texttt{find\_pivot\_kernel} & Parallel-reduction pivot search over column $k$ \\
\texttt{swap\_rows\_kernel} & Row swap in $U$, $P$, and (partially) $L$ \\
\texttt{compute\_multipliers\_kernel} & $L_{jk} = U_{jk}/U_{kk}$ for $j > k$ \\
\texttt{eliminate\_kernel} & $U_{jl} \mathrel{-}= L_{jk}U_{kl}$, with epsilon clamp \\
\texttt{apply\_permutation\_kernel} & $b' = Pb$ \\
\texttt{forward\_substitution\_kernel} & Solve $Ly = b'$ (sequential) \\
\texttt{backward\_substitution\_kernel} & Solve $Ux = y$ (sequential) \\
\bottomrule
\end{tabular}
\caption{CUDA kernels in \texttt{lu\_solve\_gpu.cu}. The elimination loop
(pivot search through elimination) runs once per pivot step
$k = 0,\dots,18$, i.e.\ each kernel above (except substitution) is launched
$\sim$19 times per program run.}
\label{tab:kernels}
\end{table}

\subsection{Build and Environment}
\begin{itemize}
  \item \textbf{GPU}: NVIDIA GeForce RTX 3090, compute capability 8.6
  \item \textbf{Compilation}: \texttt{nvcc -std=c++17 -g -O2 -lineinfo -arch=sm\_86}
    (see \texttt{Makefile}); \texttt{-lineinfo} enables nixnan's
    source-line attribution.
  \item \textbf{Instrumentation}: \texttt{LD\_PRELOAD=../../nixnan.so}
\end{itemize}

\section{Stage 1: Exp-Binade Histogram Gathering}
\label{sec:binade}

\subsection{Configuration}

\texttt{bin\_spec.json} (generated by \texttt{gen\_bin\_spec.py}) configures
nixnan's \texttt{BIN\_SPEC\_FILE} binade-monitoring feature:

\begin{lstlisting}
{
    "count": 128,
    "doublings": 7,
    "max_reports": 0,
    "f16": [],
    "bf16": [],
    "f32": [],
    "f64": [[-1022, -1019], [-1018, -1015], ..., [1022, 1023]]
}
\end{lstlisting}

\begin{itemize}
  \item \textbf{All precisions present}: \texttt{lu\_solve\_gpu.cu} computes
    only in \texttt{double}, so \texttt{f64} is the only format populated
    with exponent bins; \texttt{f16}/\texttt{bf16}/\texttt{f32} are left
    empty, which is valid per nixnan's schema and correctly reflects that
    those formats are not present in the solver's source-level arithmetic
    (see Section~\ref{sec:fp32-artifact} for a caveat about what nixnan
    itself classifies at the SASS level).
  \item \textbf{Histogram bucket size 4}: the \texttt{f64} array tiles the
    full valid double exponent range $[-1022, 1023]$ (normal exponents
    only; the reserved subnormal/inf/NaN exponent codes are excluded) into
    512 buckets of width 4.
  \item \textbf{Report threshold \texttt{count: 128}}: an initial attempt
    with \texttt{count: 1} produced roughly 13{,}700 threshold-report lines
    for a single run, because nixnan reports on \emph{every} occurrence past
    the first once \texttt{count} is 1 (see the report condition in
    Section~\ref{sec:doublings}). \texttt{128} was chosen to match nixnan's
    own convention (its auto-generated template, and the
    \texttt{SC26-expts/rd\_nixnan} example, both default to this order of
    magnitude) and produces a much more legible 99--63 reports per run here.
\end{itemize}

\subsection{Results}

\begin{table}[h]
\centering
\begin{tabular}{lrrr}
\toprule
& \texttt{matrix.csv} & \texttt{bad\_matrix.csv} \\
\midrule
Residual $\lVert Ax-b\rVert_2$ (uninstrumented) & $9.68\times10^{-16}$ & \texttt{nan} \\
Binade-threshold reports & 99 & 63 \\
FP64 NaN (repeats) & 0 (0) & 316 (2171) \\
FP64 Infinity (repeats) & 0 (0) & 4 (0) \\
FP64 Division-by-zero (repeats) & 0 (0) & 4 (4) \\
\bottomrule
\end{tabular}
\caption{Stage 1 results, from \texttt{logs/baseline/*.log} and
\texttt{logs/histogram/*.log}.}
\label{tab:stage1}
\end{table}

This reproduces the CPU tutorial's central finding on the GPU: the
well-conditioned matrix produces a clean solution and no exceptions of any
kind, while the ill-conditioned matrix triggers exactly the failure the
tutorial is built around (a near-zero pivot leading to division-by-zero,
then NaN propagation through the whole solve).

\subsection{An FP32 Artifact of Binary Instrumentation}
\label{sec:fp32-artifact}

Despite the source computing exclusively in \texttt{double}, the final
report for \texttt{bad\_matrix.csv} also shows 6 FP32 NaN detections (93
repeats). Inspecting the flagged SASS instructions
(\texttt{TOOL\_VERBOSE=1}, \texttt{PRINT\_ILL\_INSTR=1}) shows the pattern
\texttt{FFMA R8, RZ, R5, R3} in \texttt{compute\_multipliers\_kernel} and
\texttt{backward\_substitution\_kernel}. This is \texttt{ptxas}'s
``multiply-by-the-zero-register'' idiom for a register move
($0 \times R5 + R3 = R3$), used here to shuffle the 32-bit halves of a
64-bit double across the register file rather than to perform a genuine
32-bit floating-point computation. Because nixnan instruments at the SASS
opcode level, it correctly classifies \texttt{FFMA} as an \texttt{f32}
instruction and correctly notices that the multiplicand \texttt{R5} happens
to hold a 32-bit sub-word bit pattern that decodes as NaN -- but that
32-bit word is a half of a double, never interpreted as an \texttt{f32}
value by the program. This is a known class of false positive for
binary-level FP exception detectors and is worth keeping in mind when
reading ``per-format'' exception counts: a nonzero count in a format the
source code never uses does not necessarily mean that format's arithmetic
ran.

\section{Stage 2: Exponential \texttt{SAMPLING} Sweep}
\label{sec:sampling}

\subsection{Requested Schedule}

The requested instrumentation-sampling schedule was: \texttt{SAMPLING=1}
for 16 runs (``snaps''), then \texttt{2} for 16, then \texttt{4} for 16,
then \texttt{8} for 16, then reset back to \texttt{1}.

nixnan reads \texttt{SAMPLING} once at process start
(\texttt{GET\_VAR\_INT(sampling, "SAMPLING", ...)} in \texttt{src/nixnan.cu})
and there is no supported way to change it mid-run. The schedule is
therefore implemented in \texttt{run\_nixnan\_experiment.sh} as 64 separate
whole-program launches of \texttt{lu\_solve\_gpu bad\_matrix.csv} (one
``snap'' = one process), storing each run's nixnan log and stdout under
\texttt{logs/sampling/cycle\_01/sampling\_\{1,2,4,8\}/snap\_\{01..16\}}.
Since the solver is deterministic (fixed CSV input, no randomness), the 16
snaps at a fixed level are expected -- and observed -- to be byte-identical;
the stored logs faithfully record the \emph{schedule}, not run-to-run
variation.

\subsection{Finding: \texttt{SAMPLING} Has No Observable Effect}
\label{sec:sampling-bug}

\texttt{bad\_matrix.csv}'s elimination kernels (Table~\ref{tab:kernels}) are
each launched $\sim$19 times, once per pivot step, so
$\texttt{SAMPLING}\in\{1,2,4,8\}$ should visibly change how many of those
repeat invocations get instrumented. It does not: \texttt{snap\_01.log} for
\texttt{SAMPLING=1} and for \texttt{SAMPLING=8} are byte-identical
(\texttt{diff} produces no output), as are all 16 snaps within any single
level (one MD5 across all 16 files).

Reading \texttt{src/nixnan.cu} explains why. The sampling-skip check lives
in \texttt{should\_instrument()}:

\begin{lstlisting}[language=C++]
bool should_instrument(CUcontext ctx, CUfunction f) {
  std::string func_name = nvbit_get_func_name(ctx, f);   // full signature
  ...
  if (sampling != 0 && analyzed_kernels.count(func_name)) {
    if (analyzed_kernels[func_name] % sampling != 0) {
      ++analyzed_kernels[func_name];
      enable_instr = false;
    }
  }
  return enable_instr;
}
\end{lstlisting}

\noindent while the counter it reads, \texttt{analyzed\_kernels}, is
incremented in the CUDA-launch callback under a \emph{different} key:

\begin{lstlisting}[language=C++]
std::string kernel_name = nvbit_get_func_name(ctx, p->f);
std::string short_name = cut_kernel_name(kernel_name);   // truncated name
...
if (enable_instr) {
  int count = analyzed_kernels[short_name]++;            // <-- keyed on short_name
  ...
  ++analyzed_kernels[short_name];
}
\end{lstlisting}

\noindent where \texttt{cut\_kernel\_name} (\texttt{src/common.cuh})
truncates at the first \texttt{<} or \texttt{(}:

\begin{lstlisting}[language=C++]
inline std::string cut_kernel_name(const std::string& kernel_name) {
  size_t pos = kernel_name.find_first_of("<(");
  if (pos != std::string::npos) return kernel_name.substr(0, pos);
  return kernel_name;
}
\end{lstlisting}

For a kernel such as
\texttt{eliminate\_kernel(double*, double const*, int, int)}, the two keys
are \texttt{"eliminate\_kernel(double*, double const*, int, int)"} (full
signature, used in the lookup) versus \texttt{"eliminate\_kernel"}
(truncated, used for the increment). These never match for any kernel with
a non-empty parameter list -- i.e.\ essentially every real CUDA kernel --
so \texttt{analyzed\_kernels.count(func\_name)} is always \texttt{0} inside
\texttt{should\_instrument()}, the skip branch never executes, and every
invocation is instrumented irrespective of \texttt{SAMPLING}. This appears
to be a genuine bug in nixnan's \texttt{SAMPLING} feature, independent of
this experiment, and worth fixing upstream. The 64 stored logs remain a
faithful record of the requested schedule and harness, but currently
demonstrate no sampling-induced difference in coverage.

\section{A Second Documentation/Implementation Gap: \texttt{doublings}}
\label{sec:doublings}

nixnan's documentation (\texttt{README.md}, \texttt{Tutorial.md}) describes
a \texttt{doublings} field for \texttt{BIN\_SPEC\_FILE} that adaptively
scales the report threshold ($256\to512\to1024\to\cdots\to$ reset). As of
this checkout, \texttt{src/fp-histogram.cu} never reads that key. The only
threshold logic is in \texttt{record()} (\texttt{src/fp-fun.cu}):

\begin{lstlisting}[language=C++]
unsigned long long int prev = bin.increment();
if (prev % count == 0 && prev != 0) {
    // ... send report ...
}
\end{lstlisting}

which reports every time the running count is a positive multiple of the
\emph{static} \texttt{count} threshold -- with no reference to a doubling
schedule anywhere in \texttt{src/}. \texttt{bin\_spec.json} still sets
\texttt{doublings} for schema/forward compatibility, but it is currently a
no-op. This is why Stage 2's exponential schedule was implemented at the
shell-script level (Section~\ref{sec:sampling}) rather than relying on
\texttt{BIN\_SPEC\_FILE}.

\section{Conclusions}

\begin{enumerate}
  \item The GPU port of the tutorial's LU solver reproduces the CPU
    behavior exactly at the level that matters for the tutorial's purpose:
    clean convergence on a well-conditioned system, and NaN/Inf/Div-by-zero
    on an ill-conditioned one, both correctly detected and attributed by
    nixnan's exception instrumentation and its exp-binade histogram
    feature.
  \item Two tool-level gaps were found and documented with code-level
    evidence: \texttt{SAMPLING} does not filter any instrumentation in this
    build (Section~\ref{sec:sampling-bug}), and \texttt{BIN\_SPEC\_FILE}'s
    \texttt{doublings} field is unimplemented (Section~\ref{sec:doublings}).
    Neither was fixed here (out of scope for this experiment), but both are
    concrete, reproducible findings for the nixnan maintainers.
  \item Binary-level FP exception detection can misattribute a compiler
    register-move idiom as a genuine narrow-format computation
    (Section~\ref{sec:fp32-artifact}); per-format counts should be read
    alongside the source's actual precision usage, not in isolation.
\end{enumerate}

\section*{Reproduction}

\begin{lstlisting}
cd ~/repos/nixnan/SC26-expts/lu_solve_gpu
./run_nixnan_experiment.sh
\end{lstlisting}

\noindent Requires a built \texttt{../../nixnan.so} (\texttt{cd ../.. \&\& make}).
See \texttt{README.md} in this directory for file-by-file documentation and
override variables (\texttt{NIXNAN\_SO}, \texttt{CYCLES},
\texttt{SAMPLING\_MATRIX}).

\end{document}
